\documentclass{willowtreebook}
\Title{Lecture Notes on Sequences}
\Author{Noor Fatima}
\Colophon{Released for classroom use.}
\BibliographyFile{references}
\usepackage{booktabs}

\begin{document}

\chapter{Preface}
These notes use the Willow Tree Book class, a simplified derivative of
\texttt{memoir} intended for lecture notes: plain chapters, a bibliography,
and definitions and problems set apart from the running text.
\afterpreface

\chapter{Convergent Sequences}

\section{Definitions}

We call a sequence \emph{convergent}\define{convergent sequence} if it
approaches a single value as its index grows, following the standard
treatment in \cite{knuth1984texbook}.

\section{A Summary Table}

Table~\ref{tab:terms} lists the terms defined in this chapter.

\begin{table}[h]
  \centering
  \begin{tabular}{ll}
    \toprule
    Term & Introduced in \\
    \midrule
    Convergent sequence & Section~1.1 \\
    Bounded sequence     & Section~1.2 \\
    \bottomrule
  \end{tabular}
  \caption{Terms defined in this chapter.}
  \label{tab:terms}
\end{table}

\section{Bounded Sequences}

A sequence is \emph{bounded}\define{bounded sequence} if all its terms lie
within some fixed interval.

\begin{problem}{seq.bounded.convergent}
Show that every convergent sequence is bounded.
\end{problem}
\begin{answer}{seq.bounded.convergent}
If a sequence converges to \(L\), all but finitely many terms lie within
distance \(1\) of \(L\); the remaining finitely many terms are trivially
bounded, so the whole sequence is bounded.
\end{answer}

Problem~\ref{problem:seq.bounded.convergent} is a standard first exercise
\cite{knuth1984texbook}.

\chapter{Closing Remarks}

This class does not alter how a bibliography is produced: a plain
\texttt{amsplain}-styled \texttt{.bib} file, resolved the same way it would
be under \texttt{book}.

\end{document}
